\documentclass{article}
\usepackage{hyperref}
\usepackage{Style}

\nocite{*} % Comentar si quiero citar
%\addbibresource{bibliografia.bib} % Quitar el comentado si quiero usar bibliografia

\begin{document}

\begin{minipage}{2.5cm}
    \includegraphics[width=2cm]{imagen_puc.jpg}
\end{minipage}
\begin{minipage}{12cm}
    {\sc Pontificia Universidad Católica de Chile\\
    Facultad de Matemáticas\\
    Departamento de Matemática\\
    Profesor: Manuel Cabezas -- Ayudante: Benjamín Mateluna}
\end{minipage}
\vspace{2ex}

{\centerline{\bf Introducción a la Geometría - MAT1304}
\centerline{\bf Ayudantía 15}}
\centerline{\bf 30 de abril de 2026}

\vspace{1cm}
\noindent\textbf{Problema 1.} Encuentre las 3 raíces del polinomio $x^{3}+3x^{2}-2x-6$.

\vspace{5mm}
\noindent\textbf{Problema 2.} Mostrar que no existen $p,q\in\Q[x]$ tales que 
$x^{3}+101x^{2}-107x+2=p(x)q(x)$.

\vspace{5mm}
\noindent\textbf{Problema 3.} Demuestre que $\sqrt[3]{2+\sqrt{5}}+\sqrt[3]{2-\sqrt{5}}=1$.

\vspace{5mm}
\noindent\textbf{Problema 4.} Encontrar un par de números $a,b$ para que el polinomio 
$p(x)=x^{3}-ax^{2}+bx+5$ tenga solo raíces racionales y distintas.

\vspace{1cm}
\noindent\textbf{\underline{Ejercicios Propuestos:}}

\vspace{5mm}
\noindent\textbf{Extra 1.} Demuestre que $\sqrt{2}+\sqrt{3}$ es irracional.

\vspace{5mm}
\noindent\textbf{Extra 2.} Encontrar un polinomio $p(x)$ de grado $5$ tal que $p(x)=r(x)t(x)$ y
$\pm1$ son raices de $p(x)$. Donde $r(x)$ es de grado $2$ y $t(x)$ no tiene raíces racionales.

%\printbibliography % Quitar el comentado si quiero usar bibliografia

\end{document}
